
\section{Introduction}

\subsection{Motivation}

The quadratic parameter cost of transformer layers makes deployment of even 
moderately-sized models expensive. A standard BERT-base layer contains 
approximately 7 million parameters, of which 4.7 million (67\%) live in the 
four attention projections $W_Q, W_K, W_V, W_O$ and the remaining 2.4 million 
in the FFN. For DistilBERT, this drops to 3.5M and 2.4M respectively. However, 
even DistilBERT (66.9M parameters) requires 134 MB of storage at FP16, which 
is not negligible for mobile or edge deployment.

Low-rank factorization is the classical approach to reducing this cost. Given 
a linear layer $y = x W^\top + b$, one replaces $W$ with a product $W \approx U V$ 
where $U \in \mathbb{R}^{d_\text{out} \times r}$ and $V \in \mathbb{R}^{r \times d_\text{in}}$, 
$r \ll \min(d_\text{in}, d_\text{out})$. The parameter count decreases from 
$d_\text{in} d_\text{out}$ to $r(d_\text{in} + d_\text{out})$, achieving 
compression ratio $\frac{r(d_\text{in} + d_\text{out})}{d_\text{in} d_\text{out}}$. 
For a $768 \times 768$ matrix at rank $r = 64$, this is a 5.9$\times$ reduction 
(589,824 $\to$ 98,304 parameters).

Classical SVD is optimal in an unweighted Frobenius norm sense:

\[
W_r^* = \arg\min_{W_r : \text{rank}(W_r) = r} \|W - W_r\|_F
\tag{1.1}
\]

but this optimality is with respect to uniform treatment of all input 
dimensions. Transformer activations, however, have **highly non-uniform 
variance** across channels. Some input dimensions carry most of the signal; 
others are nearly constant. Reconstructing the near-constant dimensions as 
Activation-aware compression addresses this fundamental limitation by weighting the reconstruction 
objective per input dimension according to its empirical importance. The natural weight is the empirical Fisher 
diagonal, $F_i = \mathbb{E}_{\mathbf{x}} [(\nabla_{W} \mathcal{L})_{i}^2] \approx \mathbb{E}_{\mathbf{x}}[x_i^2]$, 
which measures the second-order sensitivity and variance attributable to input channel $i$. 
Reweighting the SVD input space by this diagonal quantity concentrates the low-rank approximation 
capacity precisely onto the dominant feature directions that govern downstream task representations.

Given a pretrained transformer model $\mathcal{M}$ with $N_l$ layers, each 
containing linear projections $\{W^{(\ell)}_k\}$, we seek a compressed model 
$\mathcal{M}'$ such that:

\[
\mathcal{M}' \approx \mathcal{M}, \qquad
\frac{|\theta(\mathcal{M}')|}{|\theta(\mathcal{M})|} \ll 1
\tag{1.2}
\]

where $|\theta(\cdot)|$ denotes parameter count.

We impose the following constraints:

\begin{enumerate}
\item \textbf{Strict reconstruction bound.} For each compressed layer,
  $\|W - UV\|_F / \|W\|_F \leq \tau$ for threshold $\tau = 0.40$.
\item \textbf{No retraining of base weights.} Only LoRA adapters may be trained.
\item \textbf{Memory-bounded calibration.} Activation statistics must be collected in
  $\mathcal{O}(d^2)$ memory, not $\mathcal{O}(Nd^2)$.
\item \textbf{Deterministic fallback.} If error exceeds $\tau$, the original layer is
  preserved.

\end{enumerate}
\subsection{Our Approach}

We propose \textbf{A3} (\textbf{A}ctivation-\textbf{A}ware \textbf{A}daptive factorization), 
which differs from standard SVD compression in four ways:

\begin{enumerate}
\item \textbf{Fisher-diagonal weighting.} We estimate the empirical Fisher diagonal
  $F_i = \mathbb{E}_x[x_i^2]$ from calibration data and use it to reweight
  the SVD input space: $W_\text{weighted} = W \cdot \mathrm{diag}(\sqrt{F + \epsilon})$.

\item \textbf{Streaming covariance accumulation.} We collect $F$ and $C_X$ using
  online streaming, requiring only $\mathcal{O}(d^2)$ memory regardless of
  calibration set size.

\item \textbf{Adaptive rank escalation.} If the reconstruction error at the initial
  rank $r$ exceeds $\tau$, we incrementally increase the rank up to
  $\min(d_\text{in}, d_\text{out})$, retrying the SVD until the constraint
  is met or the fallback triggers.

\item \textbf{LayerNorm recalibration.} After factorization, we rescale each
  LayerNorm's affine parameters $\gamma$ by the ratio of target to current
  activation standard deviation, absorbing distributional shift introduced
  by the factorization.

\end{enumerate}
\subsection{Contributions}

We make the following contributions:

\begin{itemize}
\item \textbf{A streaming activation-aware compression algorithm.} We propose A3,
  which combines Fisher weighting, streaming covariance accumulation, adaptive
  rank escalation, and LayerNorm recalibration into a single pipeline.
\item \textbf{A production-hardened implementation.} Our PyTorch implementation
  handles flat, nested, and causal architectures with safe fallback, strict
  Frobenius validation, and CoO(nstant) memory overhead.
\item \textbf{A comprehensive evaluation.} We benchmark on DistilBERT, RoBERTa, and
  GPT-2 on SST-2 with ablations over rank, weighting strategy, and LayerNorm
  recalibration.
\item \textbf{Full reproducibility.} Code, checkpoints, and evaluation scripts are
  released under MIT.

\end{itemize}
\subsection{Paper Organization}

Section 2 reviews related work in low-rank compression, activation-aware 
pruning, and LayerNorm recalibration. Section 3 provides background on SVD, 
the Fisher information, and streaming statistics. Section 4 presents A3. 
Section 5 describes integration with three transformer architectures. Section 
6 details the experimental setup. Section 7 reports results. Section 8 
analyzes where A3 succeeds and fails. Section 9 discusses practical guidance. 
Section 10 lists limitations. Section 11 concludes.



\subsection{Notation}
